% gcf_irrationality_schema.tex — General Irrationality Schema for Quadratic-Denominator PCFs
% Central theorem for the V_quad paper
% Revision 2: fixed d_a scope, added summary table, algorithm env, polished citations
\documentclass[11pt,a4paper]{article}
\usepackage{amsmath,amssymb,amsthm}
\usepackage{hyperref}
\usepackage{booktabs}
\usepackage{algorithmicx}
\usepackage[noend]{algpseudocode}
\usepackage{algorithm}

\newtheorem{theorem}{Theorem}
\newtheorem{lemma}[theorem]{Lemma}
\newtheorem{corollary}[theorem]{Corollary}
\newtheorem{proposition}[theorem]{Proposition}
\theoremstyle{definition}
\newtheorem{definition}[theorem]{Definition}
\theoremstyle{remark}
\newtheorem{remark}[theorem]{Remark}

\DeclareMathOperator{\lcm}{lcm}
\newcommand{\abs}[1]{\lvert #1 \rvert}

\title{General Irrationality Schema for Quadratic-Denominator\\Polynomial Continued Fractions}
\author{}
\date{\today}

\begin{document}
\maketitle

%% ═══════════════════════════════════════════════════════════════════════
\section{Setup and Notation}
%% ═══════════════════════════════════════════════════════════════════════

Consider the generalized continued fraction (GCF)
\[
  \alpha = b_0 + \cfrac{a_1}{b_1 + \cfrac{a_2}{b_2 + \cdots}}
  \eqqcolon b_0 + \mathop{K}_{n=1}^{\infty}\frac{a_n}{b_n}\,,
\]
where $a_n, b_n \in \mathbb{Q}[n]$ are polynomials in $n$ with rational coefficients.
The \emph{convergents} $p_n/q_n$ are defined by the three-term recurrence
\begin{equation}\label{eq:recurrence}
  y_n = b_n\, y_{n-1} + a_n\, y_{n-2}, \qquad n \geq 1,
\end{equation}
with initial conditions $p_{-1} = 1$, $p_0 = b_0$, $q_{-1} = 0$, $q_0 = 1$.

We write
\[
  b_n = \beta_2 n^2 + \beta_1 n + \beta_0 \quad (\beta_2 \neq 0), \qquad
  a_n = \alpha_d n^d + \cdots + \alpha_0 \quad (\alpha_d \neq 0),
\]
with $d = \deg(a_n) \eqqcolon d_a$.  We define $A_n := \prod_{k=1}^n a_k$
and $B_n := \prod_{k=1}^n b_k$ ($B_0 := 1$).


%% ═══════════════════════════════════════════════════════════════════════
\section{Poincar\'e--Perron Analysis}\label{sec:poincare}
%% ═══════════════════════════════════════════════════════════════════════

Since the coefficients of~\eqref{eq:recurrence} diverge, the classical
Poincar\'e theorem does not apply directly.  We normalize by setting
$\hat{y}_n := y_n / B_n$.  Substituting into~\eqref{eq:recurrence}
and dividing by $B_n = b_n B_{n-1}$ yields the \emph{normalized recurrence}
\begin{equation}\label{eq:normalized}
  \hat{y}_n = \hat{y}_{n-1} + c_n\, \hat{y}_{n-2},
  \qquad c_n := \frac{a_n}{b_n\, b_{n-1}}\,.
\end{equation}

Since $\deg(a_n) = d_a$ and $\deg(b_n b_{n-1}) = 4$, we have
$c_n = (\alpha_d / \beta_2^2)\, n^{d_a - 4} + O(n^{d_a - 5})$.

\medskip
\noindent\textbf{Case I: $d_a \leq 3$.}
Then $c_n \to 0$ and $\sum \abs{c_n} < \infty$ (since $d_a - 4 \leq -1$).
The limiting characteristic equation of~\eqref{eq:normalized} is
$\lambda^2 - \lambda = 0$ with roots $\lambda_+ = 1$, $\lambda_- = 0$.
By the Perron--Kreuser refinement of Poincar\'e's
theorem~\cite{perron1929,pituk2002}, the absolute summability
$\sum\abs{c_n} < \infty$ guarantees that two linearly independent
solutions of~\eqref{eq:normalized} satisfy
$\hat{y}_n^{(+)} \to L \neq 0$ and $\hat{y}_n^{(-)} \to 0$.
In original coordinates:
\begin{equation}\label{eq:growth}
  y_n^{(+)} \sim L\, B_n, \qquad y_n^{(-)} = o(B_n).
\end{equation}

\noindent\textbf{Case II: $d_a = 4$.}
Then $c_n \to c := \alpha_4/\beta_2^2$.  The characteristic equation is
$\lambda^2 - \lambda - c = 0$ with roots
$\lambda_{\pm} = (1 \pm \sqrt{1+4c}\,)/2$.
Two distinct positive roots exist iff $-\tfrac{1}{4} < c < 0$,
equivalently $-\beta_2^2/4 < \alpha_4 < 0$.
Then $y_n^{(\pm)} \sim L_{\pm}\, \lambda_{\pm}^n\, B_n$
with $0 < \lambda_- < \lambda_+ < 1$.

\medskip
\noindent\textbf{Growth of $B_n$.}
For $b_k = \beta_2 k^2 + O(k)$, Stirling's approximation gives
\begin{equation}\label{eq:Bn_growth}
  \log B_n = 2n\log n + n(\log\beta_2 - 2) + O(\log n).
\end{equation}
In particular, $B_n$ (and hence $q_n$) grows \emph{super-exponentially}.


%% ═══════════════════════════════════════════════════════════════════════
\section{Irrationality via Convergent Speed}\label{sec:irrationality}
%% ═══════════════════════════════════════════════════════════════════════

\begin{lemma}[Wronskian identity]\label{lem:wronskian}
For the recurrence~\eqref{eq:recurrence}:
\[
  W_n := p_n\, q_{n-1} - p_{n-1}\, q_n = (-1)^{n-1} A_n,
  \qquad n \geq 1.
\]
\end{lemma}

\begin{proof}
Induction on $n$.  At $n=1$: $W_1 = (b_1 b_0 + a_1)\cdot 1 - b_0 \cdot b_1 = a_1$.
Inductive step:
$W_n = a_n(p_{n-2} q_{n-1} - p_{n-1} q_{n-2}) = -a_n W_{n-1}
     = (-1)^{n-1} A_n$.
\end{proof}

\begin{lemma}[Approximation quality]\label{lem:approx}
If the GCF converges to $\alpha$, then
\[
  \alpha - \frac{p_n}{q_n}
  = \sum_{m=n+1}^{\infty} \frac{(-1)^{m-1} A_m}{q_m\, q_{m-1}}\,,
\]
and in particular
$\bigl|\alpha - p_n/q_n\bigr| = \abs{A_{n+1}} / (q_n\, q_{n+1})
 \cdot \bigl(1 + O(\abs{a_{n+2}}/\abs{b_{n+2}})\bigr)$.
\end{lemma}

\begin{proof}
From the Wronskian: $p_m/q_m - p_{m-1}/q_{m-1} = (-1)^{m-1}A_m/(q_m q_{m-1})$.
Summing from $m = n+1$ to $\infty$ and noting the ratio of successive terms
$a_{m+1} q_{m-1}/q_{m+1} = O(m^{d_a - 4}) \to 0$ for $d_a \leq 3$
gives the asymptotic.  (The factor~2 in subsequent estimates absorbs the
geometric tail for all $n$ beyond an explicit computable threshold.)
\end{proof}

\begin{theorem}[Irrationality criterion]\label{thm:irrationality}
Suppose:
\begin{enumerate}
  \item[\emph{(i)}]   $a_n \neq 0$ for all $n \geq 1$;
  \item[\emph{(ii)}]  $\abs{A_n}/\abs{q_n} \to 0$ as $n \to \infty$;
  \item[\emph{(iii)}] $p_n, q_n \in \mathbb{Z}$ for all $n \geq 0$.
\end{enumerate}
Then $\alpha$ is irrational.
\end{theorem}

\begin{proof}
Suppose $\alpha = p/q$ with $p, q \in \mathbb{Z}$, $q > 0$.
Set $\delta_n := p\,q_n - q\,p_n \in \mathbb{Z}$.

By~(i), $W_n \neq 0$, so consecutive convergents are distinct;
hence $\delta_n \neq 0$ for infinitely many $n$.
For such $n$, $\abs{\delta_n} \geq 1$, giving
\[
  \abs{\alpha - p_n/q_n} = \frac{\abs{\delta_n}}{q\,q_n} \geq \frac{1}{q\,q_n}\,.
\]
By Lemma~\ref{lem:approx}:
$\abs{\alpha - p_n/q_n} \leq 2\abs{A_{n+1}}/(q_n\, q_{n+1})$.
By~(ii), $\abs{A_{n+1}}/q_{n+1} \to 0$, so for $n$ large:
\[
  \frac{2\abs{A_{n+1}}}{q_n\, q_{n+1}}
  = \frac{2}{q_n}\cdot\frac{\abs{A_{n+1}}}{q_{n+1}}
  < \frac{1}{q\,q_n}\,.
\]
Combining: $1/(q\,q_n) \leq 2\abs{A_{n+1}}/(q_n q_{n+1}) < 1/(q\,q_n)$,
a contradiction.
\end{proof}


%% ═══════════════════════════════════════════════════════════════════════
\section{Irrationality Measure}\label{sec:mu}
%% ═══════════════════════════════════════════════════════════════════════

The following lemma makes condition~(ii) of Theorem~\ref{thm:irrationality}
concrete for polynomial GCFs.

\begin{lemma}[Denominator-growth domination]\label{lem:growth_domination}
Let $\deg(b_n) = 2$ with leading coefficient $\beta_2 > 0$
and $\deg(a_n) = d$.
\begin{enumerate}
\item[\emph{(a)}] If $d \leq 1$, then
  $\prod_{k=1}^n \abs{a_k/b_k} \to 0$.
\item[\emph{(b)}] If $d = 2$ and
  $\abs{\alpha_2} < \beta_2$, then
  $\prod_{k=1}^n \abs{a_k/b_k} \leq C\,(\abs{\alpha_2}/\beta_2)^n \to 0$.
\item[\emph{(c)}] If $d = 2$ and
  $\abs{\alpha_2} = \beta_2$, then
  $\sum_{k=1}^n \log\abs{a_k/b_k}
  = (\alpha_1/\alpha_2 - \beta_1/\beta_2) \log n + O(1)$,
  so the product vanishes iff the subleading ratio favours $b_k$.
\item[\emph{(d)}] If $d \geq 3$, then
  $\prod_{k=1}^n \abs{a_k/b_k} \to \infty$,
  so condition~\emph{(ii)} of Theorem~\ref{thm:irrationality} fails
  \textup{(}the Wronskian argument does not close\textup{)}.
\end{enumerate}
\end{lemma}

\begin{proof}
Since $q_n \sim L\, B_n$ (\S\ref{sec:poincare}, Case~I),
we have $\abs{A_n}/q_n \sim L^{-1} \prod_{k=1}^n \abs{a_k}/b_k$.
For part~(a): $\abs{a_k}/b_k = O(k^{d-2}) = O(k^{-1})$ so
$\sum \log\abs{a_k/b_k} = -\infty$.
Parts~(b) and~(c) follow from $\abs{a_k}/b_k \to \abs{\alpha_2}/\beta_2$
and the standard logarithmic estimate.
For part~(d): $\abs{a_k}/b_k \sim (\abs{\alpha_d}/\beta_2)\, k^{d-2}$
with $d-2 \geq 1$, so the product
$\prod (\abs{\alpha_d}/\beta_2) k^{d-2} \geq
(\abs{\alpha_d}/\beta_2)^n (n!)^{d-2} \to \infty$.
\end{proof}

\medskip
Table~\ref{tab:degrees} summarises the status by degree.

\begin{table}[h]
\centering
\caption{Irrationality argument by numerator degree $d = \deg(a_n)$,
  with $\deg(b_n) = 2$.}\label{tab:degrees}
\begin{tabular}{@{}cccp{6.4cm}@{}}
\toprule
$d$ & GCF converges? & $\abs{A_n}/q_n \to 0$? & Status \\
\midrule
$0$ & yes & yes (automatic) & Schema applies; $\mu=2$ \\
$1$ & yes & yes (automatic) & Schema applies; $\mu=2$ \\
$2$ & yes & iff $\abs{\alpha_2} < \beta_2$ (or borderline) &
  Schema applies under $C_4$; $\mu=2$ \\
$3$ & yes & \textbf{no} (generic) &
  Convergence holds but Wronskian argument fails.
  Irrationality requires different tools
  (e.g.\ Ap\'ery-type arithmetic cancellations). \\
$4$ & yes (eigenvalue case) & open & Deferred; see~\S\ref{sec:scope}. \\
\bottomrule
\end{tabular}
\end{table}


\begin{theorem}[Irrationality measure]\label{thm:mu}
Under the hypotheses of Theorem~\ref{thm:irrationality} with
$\deg(b_n) = 2$:
\[
  \mu(\alpha) = 2.
\]
\end{theorem}

\begin{proof}
\emph{Lower bound.}  $\mu(\alpha) \geq 2$ for any irrational real
(Dirichlet's approximation theorem).

\emph{Upper bound.}  Given $p/q \in \mathbb{Q}$ with $q > 0$, choose
$n$ such that $q_n \leq q < q_{n+1}$.  By the triangle inequality:
\[
  \abs{\alpha - p/q}
  \geq \frac{\abs{p\,q_n - q\,p_n}}{q\,q_n} - \abs{\alpha - p_n/q_n}
  \geq \frac{1}{q\,q_n} - \frac{C}{q_n\,q_{n+1}}
  = \frac{1}{q\,q_n}\Bigl(1 - \frac{Cq}{q_{n+1}}\Bigr)
  \geq \frac{1}{2q^2}
\]
for $q \leq q_{n+1}/(2C)$ (using $q_n \leq q$).

The gap ratio $q_{n+1}/q_n \sim b_{n+1} \sim \beta_2(n{+}1)^2$
is polynomial in $n$, so
\[
  \frac{\log q_{n+1}}{\log q_n}
  = 1 + \frac{2\log n + O(1)}{2n\log n + O(n)} \to 1.
\]
Hence $\mu(\alpha) \leq 1 + \limsup_{n\to\infty}
\bigl(\log q_{n+1}/\log q_n\bigr) = 2$.
\end{proof}


%% ═══════════════════════════════════════════════════════════════════════
\section{The Master Theorem}\label{sec:master}
%% ═══════════════════════════════════════════════════════════════════════

\begin{theorem}[GCF Irrationality Schema]\label{thm:master}
Let $a_n, b_n \in \mathbb{Q}[n]$ with $\deg(b_n) = 2$.
Write $b_n = \beta_2 n^2 + \beta_1 n + \beta_0$ and
$a_n = \alpha_d n^d + \cdots + \alpha_0$ with $d = \deg(a_n)$.
Suppose:
\begin{enumerate}
  \item[$(C_1)$] $\beta_2 > 0$ and $b_n > 0$ for all $n \geq 1$;
  \item[$(C_2)$] $a_n \neq 0$ for all $n \geq 1$
        \textup{(}$a_n$ has no positive integer root\textup{)};
  \item[$(C_3)$] $a_n, b_n \in \mathbb{Z}$ for all $n \geq 0$;
  \item[$(C_4)$] $\displaystyle\prod_{k=1}^n \frac{\abs{a_k}}{b_k} \to 0$
        as $n \to \infty$.
\end{enumerate}
Then:
\begin{enumerate}
  \item The GCF $\alpha := b_0 + K_{n=1}^{\infty}(a_n/b_n)$ converges.
  \item $\alpha$ is irrational.
  \item $\mu(\alpha) = 2$.
\end{enumerate}
\end{theorem}

By Lemma~\ref{lem:growth_domination}, condition $(C_4)$ is equivalent to:
\begin{itemize}
\item $d \leq 1$ \textup{(}automatic\textup{)}, or
\item $d = 2$ with $\abs{\alpha_2} < \beta_2$, or
\item $d = 2$ with $\abs{\alpha_2} = \beta_2$
      and the subleading coefficients satisfy
      $\alpha_1/\alpha_2 < \beta_1/\beta_2$
      \textup{(}negative logarithmic drift\textup{)}.
\end{itemize}
For $d \geq 3$ the product in $(C_4)$ diverges
(Lemma~\ref{lem:growth_domination}(d)) and the schema does not apply;
see Remark~\ref{rem:boundary}.

\begin{proof}
\emph{Convergence.}  Since $b_n > 0$ and $\deg(b_n) = 2$,
the normalized coefficients satisfy
$c_n = a_n / (b_n b_{n-1}) = O(n^{d-4})$.
Whenever $d \leq 3$ we have $\sum \abs{c_n} < \infty$, so convergence
follows from the Poincar\'e analysis of
\S\ref{sec:poincare}~\cite{perron1929,pituk2002}.
When $d \geq 3$ but $(C_4)$ holds by direct verification,
convergence still follows from the Seidel--Stern
criterion~\cite{lorentzen2008}: $\sum b_n / \abs{a_n}$ diverges
since $b_n / \abs{a_n} \sim (\beta_2/\abs{\alpha_d}) n^{2-d}$
and $d \leq 2$ whenever $(C_4)$ holds.
(For $d \leq 2$: each $c_n = O(n^{-2})$, giving absolute summability.)

\emph{Irrationality.} $(C_2)$ gives Theorem~\ref{thm:irrationality}(i);
$(C_3)$ gives (iii); $(C_4)$ combined with $q_n \sim L\, B_n$
(\eqref{eq:growth}) gives (ii) since
$\abs{A_n}/q_n \sim L^{-1}\prod_{k=1}^n \abs{a_k}/b_k \to 0$.
Apply Theorem~\ref{thm:irrationality}.

\emph{Irrationality measure.}
$\mu(\alpha) = 2$ follows from Theorem~\ref{thm:mu}.
\end{proof}


\begin{corollary}[Unit-numerator quadratic GCFs]\label{cor:unit}
Let $A, B, C \in \mathbb{Z}$ with $A > 0$ and $An^2 + Bn + C > 0$
for all $n \geq 1$.  Then
\[
  V(A,B,C) := 1 + \cfrac{1}{A{+}B{+}C + \cfrac{1}{4A{+}2B{+}C + \cfrac{1}{\ddots}}}
\]
is irrational with $\mu\bigl(V(A,B,C)\bigr) = 2$.
\end{corollary}

\begin{proof}
Set $a_n = 1$, $b_n = An^2 + Bn + C$, $b_0 = 1$.
Then $(C_1)$--$(C_4)$ hold trivially:
$d = 0 \leq 3$; $a_n = 1 \neq 0$;\ $a_n, b_n \in \mathbb{Z}$;
$d = 0 \leq 1$.
\end{proof}

This corollary applies to all 253 distinct constants in the standard
catalogue ($A \leq 5$, $B \in [-5,5]$, $C \leq 5$, deduplicated by value).


%% ═══════════════════════════════════════════════════════════════════════
\section{Scope and Limitations}\label{sec:scope}
%% ═══════════════════════════════════════════════════════════════════════

\begin{remark}[Boundary of the method]\label{rem:boundary}
For $\deg(a_n) = d \geq 3$ with $\deg(b_n) = 2$, the numerator product
$\abs{A_n} = \prod\abs{a_k} \sim \abs{\alpha_d}^n (n!)^d$ outgrows
$q_n \sim C\, B_n \sim C\, \beta_2^n (n!)^2$, so condition~$(C_4)$ fails.
The convergent denominators do not grow fast enough to ``absorb'' the
numerator product, and the Wronskian argument does not close.

Note, however, that the GCF may still \emph{converge} when $d = 3$
(since $\sum\abs{c_n} < \infty$), and the limit may well be irrational.
Proving irrationality in this regime requires \emph{arithmetic
cancellations}---analogous to the miracle in Ap\'ery's proof of
$\zeta(3) \notin \mathbb{Q}$~\cite{apery1979,beukers1979}, which uses
$d = 6$ with $\deg(b_n) = 3$ together with the fact that
$\lcm(1,\ldots,n)$ divides the auxiliary denominators.
Extending the schema to $d \geq 3$ is an open problem;
finding such arithmetic structure (or a substitute) for specific
two-parametric families is a natural next step.
More precisely, if one could show that $\lcm(1,\ldots,n)^s$
divides $A_n$ for some $s > 0$---so that the ``arithmetically corrected''
numerator product grows more slowly than $q_n$---the Wronskian argument
would close even for $d \geq 3$.
\end{remark}

\begin{remark}[Poincar\'e eigenvalues vs.\ Wronskian]
In the Beukers framework, irrationality measures depend on the ratio
$\log\lambda_1/\log\lambda_2$ of two Poincar\'e--Perron eigenvalues.
For the quadratic-denominator family with $d \leq 2$, the eigenvalues
are degenerate ($\lambda_+ = 1, \lambda_- = 0$), and the entire proof
runs on the Wronskian identity alone---the eigenvalue analysis is
essentially irrelevant.  Only for $d = 4$ (Case~II of \S\ref{sec:poincare})
do two distinct eigenvalues arise; this regime, which may connect to
Beukers-type constructions, is deferred to future work.
\end{remark}

\begin{remark}[Constant Wronskian as arithmetic miracle]
For unit-numerator GCFs ($a_n = 1$),
$\abs{W_n} = \abs{A_n} = 1$ for all $n$---the Wronskian is
\emph{constant} at the strongest possible value.
This plays the role of the ``arithmetic miracle'' in Beukers-type
proofs: it is trivially built into the structure rather than requiring
delicate divisibility arguments.
\end{remark}


%% ═══════════════════════════════════════════════════════════════════════
\section{Verification Algorithm}\label{sec:algorithm}
%% ═══════════════════════════════════════════════════════════════════════

Algorithm~\ref{alg:verify} implements the checkability protocol for
Theorem~\ref{thm:master}.  All conditions are decidable from the
coefficients of~$a_n$ and~$b_n$.

\begin{algorithm}[h]
\caption{\textsc{VerifyGCFIrrationality}$(a_n, b_n)$}\label{alg:verify}
\begin{algorithmic}[1]
\Require Polynomials $a_n = \sum_{i=0}^d \alpha_i n^i$,\quad
  $b_n = \beta_2 n^2 + \beta_1 n + \beta_0$
\Ensure $(\mathsf{converges},\, \mathsf{irrational},\, \mu)$ or $\mathsf{FAIL}$
\Statex
\State \textbf{--- Check $(C_1)$: Positivity ---}
\State $n^* \gets -\beta_1/(2\beta_2)$
\If{$n^* \leq 1$}
  \State \textbf{assert} $b_1 = \beta_2 + \beta_1 + \beta_0 > 0$
\Else
  \State \textbf{assert} $b_{\lfloor n^*\rfloor} > 0$ \textbf{and}
    $b_{\lceil n^*\rceil} > 0$ \textbf{and} $b_1 > 0$
\EndIf
\Statex
\State \textbf{--- Check $(C_2)$: Non-vanishing ---}
\State $R \gets \{\text{positive integer roots of } a_n\}$
  \Comment{rational root theorem}
\State \textbf{assert} $R = \varnothing$
\Statex
\State \textbf{--- Check $(C_3)$: Integrality ---}
\State \textbf{assert} $a_n(0), a_n(1), \ldots, a_n(d) \in \mathbb{Z}$
  \Comment{integer-valued polynomial test}
\State \textbf{assert} $b_n(0), b_n(1), b_n(2) \in \mathbb{Z}$
\Statex
\State \textbf{--- Check $(C_4)$: Growth domination ---}
\If{$d \leq 1$}
  \State \textbf{pass} \Comment{automatic}
\ElsIf{$d = 2$ \textbf{and} $\abs{\alpha_2} < \beta_2$}
  \State \textbf{pass} \Comment{geometric decay}
\ElsIf{$d = 2$ \textbf{and} $\abs{\alpha_2} = \beta_2$}
  \State Compute $S \gets \sum_{k=1}^{10^4} \log\abs{a_k/b_k}$
  \State \textbf{assert} $S < -10$ \Comment{numerical; flag for review}
\Else
  \State \Return \textsf{FAIL}
    \Comment{$d \geq 3$: product diverges}
\EndIf
\Statex
\State \Return $(\mathsf{True},\, \mathsf{True},\, \mu = 2)$
\end{algorithmic}
\end{algorithm}


%% ═══════════════════════════════════════════════════════════════════════
%% References
%% ═══════════════════════════════════════════════════════════════════════
\begin{thebibliography}{99}
\bibitem{perron1929}
  O.~Perron, \emph{Die Lehre von den Kettenbr\"uchen}, Band~II,
  Teubner, 1929; reprinted Chelsea, 1957.

\bibitem{pituk2002}
  M.~Pituk, More on Poincar\'e's and Perron's theorems for
  difference equations,
  \emph{J.\ Difference Equ.\ Appl.}\ \textbf{8} (2002), 201--216.

\bibitem{apery1979}
  R.~Ap\'ery, Irrationalit\'e de $\zeta(2)$ et $\zeta(3)$,
  \emph{Ast\'erisque} \textbf{61} (1979), 11--13.

\bibitem{beukers1979}
  F.~Beukers, A note on the irrationality of $\zeta(2)$ and $\zeta(3)$,
  \emph{Bull.\ London Math.\ Soc.}\ \textbf{11} (1979), 268--272.

\bibitem{sleszynski1889}
  I.~\'{S}leszy\'{n}ski, \"Uber die Convergenz der Kettenbr\"uche,
  \emph{Mat.\ Sb.}\ \textbf{14} (1889), 436--498.

\bibitem{wall1948}
  H.\,S.~Wall, \emph{Analytic Theory of Continued Fractions},
  Van Nostrand, 1948.

\bibitem{lorentzen2008}
  L.~Lorentzen and H.~Waadeland,
  \emph{Continued Fractions: Convergence Theory}, 2nd ed.,
  Atlantis Press, 2008.

\bibitem{beardon2007}
  A.~Beardon and I.~Short, The Seidel, Stern, Stolz and Van~Vleck
  theorems on continued fractions,
  \emph{Bull.\ London Math.\ Soc.}\ \textbf{42} (2010), 457--466.

\bibitem{hancl2006}
  J.~Han\v{c}l and R.~Ti\v{c}h\'{y},
  On the irrationality of certain generalized continued fractions,
  \emph{Monatsh.\ Math.}\ \textbf{148} (2006), 29--44.

\bibitem{raayoni2021}
  G.~Raayoni et al., Generating conjectures on fundamental constants
  with the Ramanujan Machine, \emph{Nature} \textbf{590} (2021), 67--73.

\bibitem{ferguson1999}
  H.\,R.\,P.~Ferguson, D.\,H.~Bailey, and S.~Arwade,
  Analysis of PSLQ, an integer relation finding algorithm,
  \emph{Math.\ Comput.}\ \textbf{68} (1999), 351--369.
\end{thebibliography}

\end{document}
